\documentclass{article}
\usepackage[a4paper,left=2.2cm,right=2.2cm,top=2.5cm,bottom=2.5cm]{geometry}

\usepackage{graphicx} % Required for inserting images
\usepackage{ctex}
\usepackage{amsmath,amssymb,amsfonts}
%添加remark环境
\usepackage{amsthm}

\newtheorem{remark}{Remark}

\title{专家信息自适应融合多属性决策方法}
\date{\today}

\begin{document}

\maketitle


多属性决策问题通常涉及专家、属性与备选方案三类基本要素，其核心任务是在吸收多源偏好信息的基础上形成方案排序结果。传统多属性决策方法一般遵循“数据收集—权重确定—方案评价”的基本流程，即首先确定参与决策的专家、评价属性和备选方案，并收集历史数据、统计信息与专家判断；随后通过赋权或聚合方法形成属性权重；最后基于某种评价规则计算各备选方案的综合得分并完成排序。然而，在多专家参与的复杂决策环境中，传统方法往往面临三方面局限：其一，对精确决策数据的依赖较强，而实际中专家更容易给出序关系、相对比较或局部判断；其二，不同专家之间的知识背景、经验水平与判断能力存在差异，但传统聚合方法通常难以刻画专家贡献的异质性；其三，即使模型能够识别“谁更重要”，也往往难以进一步刻画“比别人重要多少”，从而导致专家权重和属性权重容易出现区分不足的问题。基于此，本文提出一种专家信息自适应融合多属性决策方法（\textit{Fusion-based Adaptive Multi-attribute Decision-making}, FAM），旨在将专家提供的排序信息与相对重要性信息统一转化为可计算的权重结构，并在保证方案排序区分度的同时增强专家层和属性层权重的辨识能力。

\section{问题描述}

为刻画多专家参与下的多指标决策过程，首先定义三类基本集合。令 $\mathcal{I}$ 为专家集合，$i\in\mathcal{I}$ 表示参与群体决策的专家；令 $\mathcal{J}$ 为指标集合，$j\in\mathcal{J}$ 表示用于评价备选方案的指标维度；令 $\mathcal{K}$ 为备选方案集合，$k\in\mathcal{K}$ 表示待评价的候选方案。由此，后续专家信息的输入、最小权重单元的构建以及方案排序结果的形成，均在集合 $\mathcal{I}$、$\mathcal{J}$ 与 $\mathcal{K}$ 所构成的统一分析框架下展开。

在上述集合定义基础上，模型接收如下四类原始输入信息：
\begin{enumerate}
    \item 专家权重信息 $e_i$，表示专家 $i$ 在群体决策中的已知权重，满足
    \begin{align}
        e_i\ge 0,\qquad \forall i\in\mathcal{I}, \qquad 
        \sum_{i\in\mathcal{I}} e_i=1;
    \end{align}

    \item 指标排序信息 $p_{ij}$，表示专家 $i$ 对指标 $j$ 给出的排序，用于刻画专家对各评价指标重要性的判断；

    \item 方案排序信息 $p_{ijk}$，表示专家 $i$ 在指标 $j$ 下对方案 $k$ 给出的排序，用于反映该方案在相应指标维度下的相对优劣；

    \item 方案相对重要性比值信息 $t_{ijkk'}$，表示专家 $i$ 在指标 $j$ 下对方案 $k$ 与方案 $k'$ 的相对重要性判断，即
    \begin{align}
        w_{ijk}=t_{ijkk'}\,w_{ijk'},\qquad k,k'\in\mathcal{K}.
    \end{align}
\end{enumerate}

上述四类信息共同构成 FAM 模型的原始输入，也是后续进行数学表达与优化求解的基础。

\begin{remark}
本文对输入信息的完备性作如下区分。首先，专家权重 $e_i$ 与指标排序信息 $p_{ij}$ 被视为必要输入，即每位参与决策的专家均应能够对指标集合给出有效排序信息。因此，对任意 $i\in\mathcal{I}$，均有
\begin{align}
    \Omega_i^{\mathcal C}\neq \varnothing,
\end{align}
其中 $\Omega_i^{\mathcal C}$ 表示由专家 $i$ 的指标排序所诱导出的有效有序指标对集合。相比之下，方案排序信息 $p_{ijk}$ 与方案相对重要性比值信息 $t_{ijkk'}$ 允许部分缺失。若专家 $i$ 未在指标 $j$ 下对部分方案给出排序，则不生成相应的方案排序约束；若某组方案相对重要性比值信息未被提供，则不生成对应的比值约束。因而，后续模型中的各类约束索引集合仅由已观测到的有效输入构成，从而使模型能够兼容“指标层信息完整、方案层信息部分缺失、逐步补充”的专家输入情形。
\end{remark}

在给定专家权重 $e_i$ 的前提下，模型进一步构建最小权重单元 $w_{ijk}$，用于表示专家 $i$ 在指标 $j$ 下对方案 $k$ 的局部评价结果。基于 $w_{ijk}$，可沿不同维度进行聚合，从而得到专家内部指标权重、指标全局权重以及方案综合权重。具体定义如下：
\begin{align}
    u_{ij} &= \sum_{k\in\mathcal{K}} w_{ijk},
    && \forall i\in\mathcal{I},\ j\in\mathcal{J}, \\
    c_j &= \sum_{i\in\mathcal{I}} u_{ij}
    = \sum_{i\in\mathcal{I}}\sum_{k\in\mathcal{K}} w_{ijk},
    && \forall j\in\mathcal{J}, \\
    a_k &= \sum_{i\in\mathcal{I}}\sum_{j\in\mathcal{J}} w_{ijk},
    && \forall k\in\mathcal{K}.
\end{align}
其中，$u_{ij}$ 表示专家 $i$ 对指标 $j$ 的局部权重贡献，$c_j$ 表示指标 $j$ 的全局权重，$a_k$ 表示方案 $k$ 的综合权重。

与此同时，为保证给定专家权重能够在三维最小权重单元上得到一致表达，进一步施加如下分配约束：
\begin{align}
    \sum_{j\in\mathcal{J}}\sum_{k\in\mathcal{K}} w_{ijk}=e_i,
    \qquad \forall i\in\mathcal{I}.
\end{align}
由此，模型不再将专家权重视为待识别变量，而是将其作为外生输入，并在此基础上进一步分解得到指标层与方案层的权重结构。

为了实现由“专家判断”到“可计算权重结构”的转化，需要建立输入信息与输出变量之间的数学关系。对于排序型信息，本文采用“排序越靠前，对应权重越大”的一致性转译原则。具体而言：
\begin{enumerate}
    \item 对于指标排序信息，若在专家 $i$ 看来指标 $j$ 的排序优于指标 $j'$，则其对应局部权重贡献应不低于后者，即
    \begin{align}
        p_{ij}<p_{ij'} \Longrightarrow u_{ij}>u_{ij'};
    \end{align}

    \item 对于方案排序信息，若在专家 $i$ 和指标 $j$ 下，方案 $k$ 的排序优于方案 $k'$，则其最小权重单元应不低于后者，即
    \begin{align}
        p_{ijk}<p_{ijk'} \Longrightarrow w_{ijk}>w_{ijk'};
    \end{align}

    \item 对于专家直接给出的方案相对重要性比值信息，则有
    \begin{align}
        w_{ijk}=t_{ijkk'}\,w_{ijk'}.
    \end{align}
\end{enumerate}

需要指出的是，若仅利用上述序关系约束，则模型虽然能够识别“先后顺序”，但未必能够充分识别“差异程度”。一方面，在方案层面，若综合权重非常接近的方案被强行赋予不同名次，可能夸大细微差异；另一方面，在指标层面，若仅要求顺序一致，则同一专家内部不同指标之间的权重差异可能退化为最小可行差距，从而削弱指标权重结构的辨识能力。基于此，本文在模型中同时引入方案层级区分度变量 $d^{\mathcal A}$ 与专家内部指标区分度变量 $d_i^{\mathcal C}$，分别用于刻画方案层与指标层的最小权重差异。由此，FAM 方法不仅关注最终方案排序结果，也同时关注专家内部指标权重结构的可辨识性。

综上，$w_{ijk}$ 成为连接原始输入与最终决策结果的核心载体。通过对 $w_{ijk}$ 沿不同维度进行聚合，模型可得到专家给定权重下的指标权重 $c_j$ 以及方案综合权重 $a_k$，并最终形成备选方案的层级排序结果。至此，专家原始判断信息便完成了由“主观输入”向“可约束、可计算、可解释的权重结构”的映射。


\section{模型建立}

\subsection{模型构建思路}

在上述问题描述基础上，本文进一步构建 FAM 的数学规划模型。模型的核心思想是：在满足给定专家权重、指标排序信息、方案排序信息以及方案相对重要性比值信息的前提下，基于统一的最小权重单元 $w_{ijk}$，同时生成方案综合权重与指标层权重结构。考虑到实际决策中，部分方案的综合表现可能非常接近，若仍强行赋予不同名次，可能夸大细微差异，因此本文允许综合权重相同的方案处于同一排名层级。与此同时，为避免同一专家内部不同指标的权重塌缩为近似相等的结果，模型进一步引入指标层区分度变量，刻画专家内部指标重要性差异的辨识能力。

因此，本文不再采用分层词典序优化，而是构建统一的加权单目标模型，在同一优化框架下同时提升方案层级排序结果的辨识性与专家内部指标权重结构的辨识性。设 $d^{\mathcal A}$ 表示不同方案层级之间的最小区分度，$d_i^{\mathcal C}$ 表示专家 $i$ 内部不同指标之间的最小区分度，则模型目标写为
\begin{align}
    \max \quad \alpha d^{\mathcal A}+\beta \sum_{i\in\mathcal I} d_i^{\mathcal C},
\end{align}
其中，$\alpha>0$、$\beta>0$ 为目标权重参数，通常取 $\alpha>\beta$，以体现方案层级排序结果相较于指标层结构辨识性的优先性。由于在当前设定下
\begin{align}
    \sum_{i\in\mathcal I} e_i=1,
\end{align}
且各区分度变量均由最小权重单元的差值所界定，因此 $d^{\mathcal A}$ 与 $\sum_{i\in\mathcal I} d_i^{\mathcal C}$ 处于可比的尺度范围内，从而可以在统一目标函数中进行加权处理。

\subsection{方案层与指标层的统一建模}

设方案 $k$ 的综合权重为
\begin{align}
    a_k=\sum_{i\in\mathcal{I}}\sum_{j\in\mathcal{J}} w_{ijk},
    \qquad \forall k\in\mathcal{K}.
\end{align}
同时，将专家内部指标权重与指标全局权重写为
\begin{align}
    u_{ij} &= \sum_{k\in\mathcal{K}} w_{ijk},
    \qquad \forall i\in\mathcal{I},\ j\in\mathcal{J}, \\
    c_j &= \sum_{i\in\mathcal{I}} u_{ij},
    \qquad \forall j\in\mathcal{J}.
\end{align}
此外，为与已知专家权重保持一致，要求
\begin{align}
    \sum_{j\in\mathcal{J}}\sum_{k\in\mathcal{K}} w_{ijk}=e_i,
    \qquad \forall i\in\mathcal{I}.
\end{align}
上述关系表明，专家层、指标层与方案层均由同一组最小权重单元 $w_{ijk}$ 统一生成，从而保证模型结构的一致性。

为了提高方案层级排序结果的辨识性，同时允许综合表现非常接近的方案处于同一层级，设 $\tilde a_r$ 表示按综合权重排序后位于第 $r$ 位的方案权重，$r=1,\dots,|\mathcal K|$。进一步引入二元变量
\begin{align}
    y_r=
    \begin{cases}
        1, & \text{若第 } r \text{ 位与第 } r+1 \text{ 位属于不同层级},\\
        0, & \text{若第 } r \text{ 位与第 } r+1 \text{ 位属于同一层级},
    \end{cases}
    \qquad r=1,\dots,|\mathcal K|-1.
\end{align}
由此，$d^{\mathcal A}$ 不再刻画所有相邻排序位置之间的最小差距，而是仅刻画不同方案层级之间的最小差距。当相邻两个排序位置被判定为同层时，模型允许其综合权重完全相等；当二者被判定为不同层级时，则要求其权重差异至少达到给定阈值，并进一步参与方案层区分度 $d^{\mathcal A}$ 的刻画。

另一方面，为增强模型对专家内部指标权重结构的刻画能力，本文进一步引入专家内部指标区分度变量
\begin{align}
    d_i^{\mathcal C}\ge 0,
    \qquad \forall i\in\mathcal I,
\end{align}
用于描述专家 $i$ 内部不同指标之间的最小权重差异。显然，$d_i^{\mathcal C}$ 越大，说明专家 $i$ 对不同指标的重要性区分越明显。

\subsection{排序一致性与比值一致性约束}

对于指标排序信息与方案排序信息，分别建立顺序一致性约束。考虑到严格不等式在数学规划中不便直接处理，本文引入足够小的正常数 $\varepsilon>0$，并将排序一致性约束写为
\begin{align}
    u_{ij}-u_{ij'} &\ge \varepsilon,
    && \forall i\in\mathcal I,\ \forall (j,j')\in\Omega_i^{\mathcal C}, \\
    w_{ijk}-w_{ijk'} &\ge \varepsilon,
    && \forall i\in\mathcal I,\ \forall j\in\mathcal J,\ \forall (k,k')\in\Omega_{ij}^{\mathcal A},
\end{align}
其中，
\begin{align}
    \Omega_i^{\mathcal C}
    := \{(j,j')\mid p_{ij}<p_{ij'}\},
    \qquad \forall i\in\mathcal I,
\end{align}
表示由专家 $i$ 的指标排序所诱导出的有效有序指标对集合；
\begin{align}
    \Omega_{ij}^{\mathcal A}
    := \{(k,k')\mid p_{ijk}<p_{ijk'}\},
    \qquad \forall i\in\mathcal I,\ \forall j\in\mathcal J,
\end{align}
表示专家 $i$ 在指标 $j$ 下由方案排序所诱导出的有效有序方案对集合。

在此基础上，为进一步刻画专家内部指标权重之间的最小区分度，对任意有效指标排序对 $(j,j')\in\Omega_i^{\mathcal C}$，进一步施加
\begin{align}
    u_{ij}-u_{ij'} \ge d_i^{\mathcal C},
    \qquad \forall i\in\mathcal I,\ \forall (j,j')\in\Omega_i^{\mathcal C}.
\end{align}
该约束表明：对专家 $i$ 而言，其排序更靠前的指标相对于排序更靠后的指标，不仅应满足最小容差 $\varepsilon$，还应共同支撑专家内部指标区分度变量 $d_i^{\mathcal C}$。由此，$\varepsilon$ 用于保证顺序一致性的严格性，而 $d_i^{\mathcal C}$ 则用于刻画并提升指标权重结构的辨识能力。

对于专家直接给出的方案相对重要性比值信息，建立如下约束：
\begin{align}
    w_{ijk}=t_{ijkk'}\,w_{ijk'},
    \qquad \forall (i,j,k,k')\in\Omega^{\mathcal R},
\end{align}
其中
\begin{align}
    \Omega^{\mathcal R}
    := \{(i,j,k,k')\mid \text{专家 } i \text{ 已在指标 } j \text{ 下给出方案 } k \text{ 与 } k' \text{ 的相对重要性比值}\}.
\end{align}

此外，为保证权重变量的可比性与可解释性，引入非负约束：
\begin{align}
    w_{ijk}\ge 0,
    \qquad \forall i\in\mathcal I,\ j\in\mathcal J,\ k\in\mathcal K.
\end{align}
由于已知专家权重满足 $\sum_{i\in\mathcal I} e_i=1$，再结合
\begin{align*}
    \sum_{j\in\mathcal J}\sum_{k\in\mathcal K} w_{ijk}=e_i,
    \qquad \forall i\in\mathcal I,
\end{align*}
可自动推出
\begin{align*}
    \sum_{i\in\mathcal I}\sum_{j\in\mathcal J}\sum_{k\in\mathcal K} w_{ijk}=1.
\end{align*}
因此，在当前设定下无需再额外施加总体归一化约束。

\subsection{允许同层的方案层级线性化表达}

由于目标函数中涉及按综合权重排序后的方案位置权重 $\tilde a_r$，其本质上隐含了方案排序结构，难以直接写为标准数学规划形式。为此，进一步引入排序位置权重变量 $\tilde a_r$、$0$--$1$ 位置分配变量 $x_{kr}$，其中
\begin{align}
    x_{kr}=
    \begin{cases}
        1, & \text{若方案 } k \text{ 位于第 } r \text{ 位},\\
        0, & \text{否则},
    \end{cases}
    \qquad \forall k\in\mathcal K,\ r=1,\dots,|\mathcal K|.
\end{align}
这里，$\tilde a_r$ 表示排序后位于第 $r$ 位的方案综合权重。于是，方案综合权重与排序位置之间可通过如下约束进行线性化联结：
\begin{align}
    \sum_{r=1}^{|\mathcal K|} x_{kr}=1,
    && \forall k\in\mathcal K, \\
    \sum_{k\in\mathcal K} x_{kr}=1,
    && \forall r=1,\dots,|\mathcal K|, \\
    a_k-\tilde a_r \le M(1-x_{kr}),
    && \forall k\in\mathcal K,\ \forall r=1,\dots,|\mathcal K|, \\
    \tilde a_r-a_k \le M(1-x_{kr}),
    && \forall k\in\mathcal K,\ \forall r=1,\dots,|\mathcal K|,
\end{align}
其中，$M$ 为足够大的正常数。

与此同时，为表示排序后方案权重的非增结构，施加
\begin{align}
    \tilde a_r\ge \tilde a_{r+1},
    \qquad \forall r=1,\dots,|\mathcal K|-1.
\end{align}

进一步地，为刻画相邻位置是否属于同一层级，给定层级判定阈值 $\phi>0$，并建立如下约束：
\begin{align}
    \tilde a_r-\tilde a_{r+1} \ge \phi y_r,
    && \forall r=1,\dots,|\mathcal K|-1, \\
    \tilde a_r-\tilde a_{r+1} \le M y_r,
    && \forall r=1,\dots,|\mathcal K|-1.
\end{align}
上述约束共同刻画了层级判定逻辑：当 $y_r=0$ 时，有
\begin{align*}
    \tilde a_r-\tilde a_{r+1}=0,
\end{align*}
即相邻两个位置属于同一层级并具有相等权重；当 $y_r=1$ 时，有
\begin{align*}
    \tilde a_r-\tilde a_{r+1}\ge \phi,
\end{align*}
即二者属于不同层级并至少保持 $\phi$ 的最小间隔。

在此基础上，为使 $d^{\mathcal A}$ 刻画不同方案层级之间的最小区分度，进一步施加
\begin{align}
    d^{\mathcal A}\le \lambda_r(\tilde a_r-\tilde a_{r+1}) + M(1-y_r),
    \qquad \forall r=1,\dots,|\mathcal K|-1.
\end{align}
其中，$\lambda_r>0$ 为排序位置 $r$ 对应的差异调节参数。当 $y_r=1$ 时，上式退化为
\begin{align*}
    d^{\mathcal A}\le \lambda_r(\tilde a_r-\tilde a_{r+1}),
\end{align*}
表明仅有不同层级之间的相邻差值参与 $d^{\mathcal A}$ 的刻画；当 $y_r=0$ 时，该约束自动放松，不再要求这一对相邻位置为方案层区分度提供支撑。若无特别偏好，实际应用中可取
\begin{align}
    \lambda_r=1,
    \qquad \forall r=1,\dots,|\mathcal K|-1.
\end{align}

\subsection{专家信息自适应融合多属性决策模型}

综上，结合方案层区分、专家内部指标区分以及层级划分控制，本文构建的 FAM 统一加权优化模型可表述为如下混合整数线性规划：
\begin{align}
    \max \quad 
    & \alpha d^{\mathcal A}
    + \beta \sum_{i\in\mathcal I} d_i^{\mathcal C}
    + \eta \sum_{r=1}^{|\mathcal K|-1} y_r \\
    \text{s.t.}\quad
    & a_k=\sum_{i\in\mathcal I}\sum_{j\in\mathcal J} w_{ijk},
    && \forall k\in\mathcal K, \\
    & u_{ij}=\sum_{k\in\mathcal K} w_{ijk},
    && \forall i\in\mathcal I,\ j\in\mathcal J, \\
    & c_j=\sum_{i\in\mathcal I} u_{ij},
    && \forall j\in\mathcal J, \\
    & \sum_{j\in\mathcal J}\sum_{k\in\mathcal K} w_{ijk}=e_i,
    && \forall i\in\mathcal I, \\
    & \sum_{r=1}^{|\mathcal K|} x_{kr}=1,
    && \forall k\in\mathcal K, \\
    & \sum_{k\in\mathcal K} x_{kr}=1,
    && \forall r=1,\dots,|\mathcal K|, \\
    & a_k-\tilde a_r \le M(1-x_{kr}),
    && \forall k\in\mathcal K,\ \forall r=1,\dots,|\mathcal K|, \\
    & \tilde a_r-a_k \le M(1-x_{kr}),
    && \forall k\in\mathcal K,\ \forall r=1,\dots,|\mathcal K|, \\
    & \tilde a_r\ge \tilde a_{r+1},
    && \forall r=1,\dots,|\mathcal K|-1, \\
    & \tilde a_r-\tilde a_{r+1} \ge \phi y_r,
    && \forall r=1,\dots,|\mathcal K|-1, \\
    & \tilde a_r-\tilde a_{r+1} \le M y_r,
    && \forall r=1,\dots,|\mathcal K|-1, \\
    & d^{\mathcal A}\le \lambda_r(\tilde a_r-\tilde a_{r+1}) + M(1-y_r),
    && \forall r=1,\dots,|\mathcal K|-1, \\
    & d^{\mathcal A}\le \bar M \sum_{r=1}^{|\mathcal K|-1} y_r, \\
    & \sum_{r=1}^{|\mathcal K|-1} y_r \ge \underline L, \\
    & u_{ij}-u_{ij'} \ge \varepsilon,
    && \forall i\in\mathcal I,\ \forall (j,j')\in\Omega_i^{\mathcal C}, \\
    & u_{ij}-u_{ij'} \ge d_i^{\mathcal C},
    && \forall i\in\mathcal I,\ \forall (j,j')\in\Omega_i^{\mathcal C}, \\
    & w_{ijk}-w_{ijk'} \ge \varepsilon,
    && \forall i\in\mathcal I,\ \forall j\in\mathcal J,\ \forall (k,k')\in\Omega_{ij}^{\mathcal A}, \\
    & w_{ijk}=t_{ijkk'}\,w_{ijk'},
    && \forall (i,j,k,k')\in\Omega^{\mathcal R}, \\
    & w_{ijk}\ge 0,
    && \forall i\in\mathcal I,\ j\in\mathcal J,\ k\in\mathcal K, \\
    & d^{\mathcal A}\ge 0, \\
    & d_i^{\mathcal C}\ge 0,
    && \forall i\in\mathcal I, \\
    & \tilde a_r\ge 0,
    && \forall r=1,\dots,|\mathcal K|, \\
    & x_{kr}\in\{0,1\},
    && \forall k\in\mathcal K,\ \forall r=1,\dots,|\mathcal K|, \\
    & y_r\in\{0,1\},
    && \forall r=1,\dots,|\mathcal K|-1.
\end{align}

上述模型中，目标函数的第一项用于提高不同方案层级之间的最小区分度，第二项用于提高各专家内部指标权重之间的最小区分度，第三项则用于在满足前两类区分要求的基础上，适度鼓励模型形成更清晰的层级划分结果。参数 $\alpha$ 与 $\beta$ 分别反映决策者对“方案层级辨识性”和“指标结构辨识性”的相对重视程度，而参数 $\eta$ 为一较小正数，用于避免模型在无必要时过度合并相邻方案层级。

约束
\begin{align}
    \tilde a_r-\tilde a_{r+1} \ge \phi y_r,\qquad
    \tilde a_r-\tilde a_{r+1} \le M y_r
\end{align}
共同刻画了相邻排序位置是否属于不同层级的逻辑关系：当 $y_r=0$ 时，相邻位置被判定为同一层级，且其综合权重相等；当 $y_r=1$ 时，相邻位置属于不同层级，且其综合权重差异至少达到阈值 $\phi$。在此基础上，
\begin{align}
    d^{\mathcal A}\le \lambda_r(\tilde a_r-\tilde a_{r+1}) + M(1-y_r)
\end{align}
保证仅有被判定为不同层级的相邻位置差异对方案层区分度 $d^{\mathcal A}$ 提供有效支撑。

进一步地，为避免全部方案被判定为同层时 $d^{\mathcal A}$ 被空放松，本文额外引入层级保护约束
\begin{align}
    d^{\mathcal A}\le \bar M \sum_{r=1}^{|\mathcal K|-1} y_r,
\end{align}
其中，$\bar M$ 为足够大的正常数。该约束表明：若不存在任何有效层级间隔，即 $\sum_r y_r=0$，则必有 $d^{\mathcal A}=0$。因此，方案层区分度必须建立在实际的层级分离基础之上。此外，约束
\begin{align}
    \sum_{r=1}^{|\mathcal K|-1} y_r \ge \underline L
\end{align}
用于控制最少分层数，其中 $\underline L$ 为决策者预先给定的最小层级间隔数量；当不需要强制规定最少分层数时，可取 $\underline L=0$。

相较于原先仅通过单纯区分度刻画方案排序质量的模型，上述修订后的统一加权目标不仅保留了方案层与指标层的联合优化结构，而且进一步增强了层级划分结果的可解释性与稳定性。也就是说，模型不再仅追求“存在差异”，而是进一步追求“差异能够体现在有效层级分离上”。这种处理使模型更能避免无意义的同层判定现象，也更便于与实际求解器实现及前端交互界面相衔接。
\subsection{求解结果的聚合与输出}

在完成统一优化求解后，模型首先得到最小权重单元 $w_{ijk}^{\star}$、方案位置分配变量 $x_{kr}^{\star}$、排序位置权重变量 $\tilde a_r^{\star}$、层级判定变量 $y_r^{\star}$，以及区分度变量 $d^{\mathcal A\star}$ 和 $d_i^{\mathcal C\star}$。在此基础上，可进一步沿不同维度进行聚合并形成最终输出结果。

专家 $i$ 对指标 $j$ 的局部权重贡献为
\begin{align}
    u_{ij}^{\star}=\sum_{k\in\mathcal K} w_{ijk}^{\star},
    \qquad \forall i\in\mathcal I,\ j\in\mathcal J.
\end{align}

指标 $j$ 的全局权重为
\begin{align}
    c_j^{\star}=\sum_{i\in\mathcal I}\sum_{k\in\mathcal K} w_{ijk}^{\star},
    \qquad \forall j\in\mathcal J.
\end{align}

方案 $k$ 的综合权重为
\begin{align}
    a_k^{\star}=\sum_{i\in\mathcal I}\sum_{j\in\mathcal J} w_{ijk}^{\star},
    \qquad \forall k\in\mathcal K.
\end{align}

此外，由于变量 $x_{kr}^{\star}$ 已建立方案与排序位置之间的一一对应关系，因此方案 $k$ 的最终排序位置可写为
\begin{align}
    R_k^{\star}=\sum_{r=1}^{|\mathcal K|} r\,x_{kr}^{\star},
    \qquad \forall k\in\mathcal K,
\end{align}
其中，$R_k^{\star}$ 越小，表示方案 $k$ 的排序越靠前。

进一步地，根据层级判定变量 $y_r^\star$ 的取值，可识别最终排序中的并列层级结构。若某一相邻位置满足 $y_r^\star=0$，则有
\begin{align*}
    \tilde a_r^\star=\tilde a_{r+1}^\star,
\end{align*}
表明对应方案属于同一排名层级；若 $y_r^\star=1$，则表明二者属于不同层级。由此，模型不仅能够输出方案的线性排序位置，还能够进一步识别允许并列情形下的层级排序结果。

需要说明的是，在当前模型设定下，专家权重 $e_i$ 为外生给定参数，因此其不再作为模型求解输出，而是作为约束条件贯穿于最小权重单元的分配过程中。故模型最终输出主要包括：各专家对指标的局部权重贡献、各指标的全局权重、各方案的综合权重及其排序结果、方案层级划分结果，以及方案层与指标层的区分度指标。
\section{与 OPA 模型的关系说明}

为说明本文模型与经典 OPA 模型之间的关系，有必要先给出 OPA 的基本形式。按照当前 PPT 中采用的记号，OPA 以连续排序方案之间的权重差异以及末位方案权重为核心，构建如下最大最小优化模型：
\begin{align}
    \max \min & \left\{ijr\bigl(W_{ijk}^{r}-W_{ijk}^{r+1}\bigr),\ ijmW_{ijk}^{m}\right\}\\
  \text{s.t.}  & \sum_{i=1}^{p}\sum_{j=1}^{n}\sum_{k=1}^{m} W_{ijk} =1, \\
  &  W_{ijk}  \ge 0, \qquad \forall i,j,k.
\end{align}
其中，$W_{ijk}^{r}$ 表示在专家 $i$、属性 $j$ 条件下，排序位置为 $r$ 的方案所对应的连续权重，$W_{ijk}^{r+1}$ 表示其后一排序位置的连续权重。该模型的目标在于：一方面尽可能拉大连续排序方案之间的权重差异，另一方面保证末位方案仍具有正的权重，从而避免权重完全退化到单点极值。

由于上述模型属于典型的最大最小形式，通常可通过引入辅助变量 $Z$ 将其转化为线性规划形式，即
\begin{align}
    \max \quad & Z \\
    \text{s.t.}\quad 
    & Z \le ijr\bigl(W_{ijk}^{r}-W_{ijk}^{r+1}\bigr), && \forall i,j,k,r, \\
    & Z \le ijmW_{ijk}^{m}, && \forall i,j,k, \\
    & \sum_{i=1}^{p}\sum_{j=1}^{n}\sum_{k=1}^{m} W_{ijk}=1, & & \\
    & W_{ijk}\ge 0, && \forall i,j,k.
\end{align}
其中，$Z$ 表示上述两类量中的最小值。通过最大化 $Z$，即可将原始的最大最小问题转化为标准线性规划问题，从而便于求解。

从建模思想上看，本文模型与 OPA 模型具有明显的相似性。首先，二者都以排序信息为基础输入，并试图将专家偏好转化为数学规划中的约束与目标。其次，二者都以权重变量为核心求解对象，并通过非负约束与归一化约束保证结果具有可比性与可解释性。再次，二者都强调排序结果之间的区分度，即希望排序更靠前的对象在权重上体现出相应优势，而不是仅给出一个不可区分的平坦分布。

但需要说明的是，本文模型与 OPA 并不是严格等价关系。OPA 的目标函数直接作用于局部连续排序权重 $W_{ijk}^{r}$，其关注重点是局部排序位置之间的权重差异；而本文模型则以 $w_{ijk}$ 作为最小局部权重单元，并进一步通过聚合得到方案综合权重 $sw_k$，其目标函数最终作用于方案综合权重之间的差异。因此，与其说 OPA 是本文模型的严格简化形式，不如说它是本文模型在仅保留排序输入、忽略比值输入、并仅针对局部连续权重进行优化时的一种参照形式。换言之，本文模型在保留 OPA 基本排序建模思想的同时，进一步扩展到了专家—属性—方案三维结构以及相对重要性比值输入的更一般情形。

\subsection{OPA 线性化模型的进一步说明}

在将 OPA 的最大最小目标转化为线性规划形式之后，其模型结构可以作进一步解释。具体而言，约束
\begin{align}
    Z\le ijr\bigl(W_{ijk}^{r}-W_{ijk}^{r+1}\bigr)
\end{align}
体现了连续排序方案之间权重差异的控制思想，即希望排序相邻的方案在权重上保持尽可能大的区分度；约束
\begin{align}
    Z\le ijmW_{ijk}^{m}
\end{align}
则用于保证末位方案仍具有一定权重，从而避免其权重退化为零。与此同时，归一化约束
\begin{align}
    \sum_{i=1}^{p}\sum_{j=1}^{n}\sum_{k=1}^{m}W_{ijk}=1
\end{align}
与非负约束
\begin{align}
    W_{ijk}\ge 0
\end{align}
共同保证了权重结果的可比性与可解释性。

需要进一步指出的是，该线性化模型的可行域中存在一类极端退化解，即仅有一个权重变量取值为 $1$，而其余 $\left(m\times n\times p-1\right)$ 个权重变量均取值为 $0$。此类解满足非负约束与归一化约束，因此从约束意义上看属于可行解。然而，这类解并不意味着必然构成模型的最优解。其原因在于，OPA 的目标仍然要求尽可能拉开连续排序方案之间的权重差异，并同时保证末位方案具有一定权重；而极端退化解往往会导致权重分布过于集中，难以体现排序结构下连续对象之间的区分关系。因此，这种退化情形更适合被理解为模型可行域中的边界现象，而不应直接视为模型所追求的理想最优结果。

从这一意义上看，OPA 的线性化模型一方面保留了排序驱动的权重生成思想，另一方面也暴露出其在局部权重空间中的边界可行现象。这一点对于理解本文模型与 OPA 的关系是有帮助的：OPA 提供了一个以排序差异为核心的经典建模框架，而本文模型则在此基础上进一步引入专家、属性与方案三维结构以及相对重要性比值输入，从而将局部排序权重的表达扩展到更一般的专家信息融合场景之中。
\end{document}
